\section{Introduction}
\label{sec:intro}

Continuous physiological monitoring at the edge, a wearable or implanted device
classifying an electrocardiogram beat by beat, runs on a joule budget with no wall
socket behind it. That budget is spent less on arithmetic than on moving data to
where the arithmetic happens. Horowitz puts an off-chip DRAM access at about
\SI{10}{\pico\joule} per bit, \SI{0.6}{\nano\joule} for eight bytes, and notes that
even with improved I/O this stays large compared with a core processor operation,
because requests and data still have to cross the processor and memory chips
\cite{Horowitz2014_ISSCC}. Sze et al.\ state the consequence for neural networks:
DRAM accesses require up to several orders of magnitude more energy than
computation, so for deep network inference the bottleneck is memory access rather
than the arithmetic \cite{Sze2017_ProcIEEE}. Purpose-built digital accelerators such
as Eyeriss reduce this cost by keeping data resident in local register files and
reusing it across many operations \cite{Chen2016_JSSC}, but the asymmetry between a
compute operation and a memory access does not disappear; it is amortized, not
removed.

Analog in-memory computing removes most of that movement by performing the
multiply-accumulate inside the memory array itself. A resistive-memory macro reaches
\num{65.2} to \num{74.0}~TOPS/W depending on precision \cite{Wan2022_Nature}; a
64-core phase-change-memory chip reaches \num{63.1}~TOPS of peak throughput at
\num{9.76}~TOPS/W in its low-precision read mode \cite{LeGallo2023_NatElectron}; and
a multi-level-cell ferroelectric crossbar in 28~nm reaches \num{885.4}~TOPS/W
\cite{Soliman2023_NatComm}. None of that is free. Analog weights carry device-level
nonidealities that a digital cell does not, and the circuitry needed to read an
analog level back out competes with the array for the energy budget: in one measured
compute-in-memory macro the analog-to-digital converter alone occupies about
\SI{90}{\percent} of the computing time \cite{Fu2024_JSSC}, and the
crossbar-plus-converter power in the ferroelectric macro above is measured, not
modelled, at \SI{153.6}{\micro\watt} \cite{Soliman2023_NatComm}. This thesis's own
accounting in Section~\ref{sec:energy} lands on the same problem from the device
side: the computing core of the deployed classifier costs \SI{535.8}{\pico\joule}
per heartbeat, and the analog-to-digital conversion a complete chip would need,
under a standard \SI{1}{\pico\joule}-per-conversion assumption, costs \num{342}
times that. How useful an analog synapse is depends on the system built around it,
and that system's cost is set upstream of the device itself. This thesis states that
before reporting the device rather than after, because the device number on its own
does not support a system claim.

A ferroelectric field-effect transistor (FeFET) stores its state as remanent
polarization in a hafnium-zirconium-oxide (\ce{Hf}\ce{Zr}\ce{O_x}, HZO) gate
dielectric rather than as trapped charge. A train of gate pulses can set its channel
conductance to one of many stable levels, and that level is read non-destructively
at zero or near-zero gate bias. Published FeFET analog synapses report five-bit
precision with a $45\times$ tunable conductance range \cite{Jerry2017_IEDM} and,
with a superlattice gate stack designed to linearize the weight update, seven-bit
precision at \SI{94.1}{\percent} accuracy on an online-trained MNIST task
\cite{Aabrar2022_TED}. Gate-all-around (nanosheet) geometry is the industry's answer
to the electrostatic control a scaled channel needs. Wrapping the gate around all
faces of the channel rather than three gives lower subthreshold swing, tighter
drain-induced barrier lowering and better short-channel immunity than a FinFET at
equal effective width
\cite{Panigrahy2024_IEEEAccess,Sreenivasulu2024_IEEEAccess}. Pairing a ferroelectric
gate stack with that geometry is a reasonable engineering choice on both counts: the
ferroelectric supplies non-volatile analog state, and the gate-all-around channel
supplies the electrostatic control needed to read that state cleanly at a scaled
node.

That pairing is not new, and this thesis does not claim it is. Stacked-nanosheet
gate-all-around FeFETs have been fabricated: a germanium-silicon-channel device
reaching a \SI{1.8}{\volt} memory window at a \SI{2}{\volt} write voltage and
endurance beyond $10^{11}$ cycles \cite{Chen2023_VLSI}, an interfacial-layer-free
follow-on of the same device family \cite{Hsieh2024_TED}, and a double-HZO
three-dimensional gate-all-around nanosheet device built specifically to homogenize
the corner field \cite{Liao2022_VLSI}. Gate-all-around nanowire cross-sections with a
ferroelectric gate stack have also been demonstrated, both silicon-compatible
\cite{Lee2021_JEDS} and on a vertical III-V channel \cite{Persson2022_EDL}. Closest to
the device studied here, Lin and Su modelled a stacked-nanosheet FeFET explicitly as
a synapse: a Monte Carlo study, under a nucleation-limited-switching ferroelectric
model and under cycle-to-cycle variation, of the same device concept and the same
application this thesis targets \cite{Lin2024_OJNano,Lin2023_SNW}. A ferroelectric
gate-all-around device, and even a ferroelectric gate-all-around synapse, is
therefore already published work.

What Lin and Su's published study does not report (and, so far as this review found,
neither does the wider fabricated-device literature cited above) is a device
calibrated against a measured transfer characteristic, a measured write-energy
budget, a measured cycle-to-cycle repeat statistic, or a trained network deployed
onto one. This thesis supplies those for one device. Section~\ref{sec:method}
calibrates a Preisach ferroelectric model against published measurements of a
comparable gate-all-around FeFET: the memory window, both threshold voltages, the
subthreshold slope and the turn-on shape are fitted; the absolute current level is
not, and Section~\ref{sec:norm} states that boundary rather than leaving it implicit.
The model carries no wake-up, fatigue or imprint term, so no cycling-endurance claim
is made anywhere in this thesis; Section~\ref{sec:limits} states that plainly as
well. Section~\ref{sec:device} measures a fifteen-level potentiation ladder at about
\SI{9}{\atto\joule} per switching pulse, with eight of those levels separable
open-loop under the device's own measured cycle-to-cycle noise and roughly sixteen
achievable across the same range under write-verify.

Section~\ref{sec:network} carries that ladder through to three results that this
thesis takes to be its substance rather than an illustration of it. First, where the
levels sit matters more than how many there are: at a fixed count of eight, four
different placements of the interior levels gave accuracies from \num{0.336} to
\num{0.491}, while spacing the same eight levels evenly gave \num{0.804}
(Section~\ref{sec:placement}). Second, level placement on this device is set by
fixed interface charge rather than by ferroelectric thickness: a \SI{10}{\percent}
perturbation of interface charge moves levels by \num{1.31} decades, while a
\SI{0.3}{\nano\metre} perturbation of the ferroelectric thickness moves them by only
\num{0.07}. That is the reverse of the ordering a thickness-first intuition would
predict, and it is the most directly actionable number in this thesis for a process
engineer (Section~\ref{sec:interface}). Third, being able to tell levels apart is
not the same as the network working: three independent measures of readout
separability all fail to track deployed accuracy, while the deployed weight's own
placement error does, with a Spearman correlation of \num{-0.85}.

The network is a spiking recurrent classifier for single-lead electrocardiogram
beats. Its architecture and MIT-BIH task originate with Yuan et al.'s
memristor-based system, which reports \SI{95.83}{\percent} on the same four-class
problem with its own software synapses \cite{Yuan2023_NatComm}. Deployed here on the
measured fifteen-level ladder by post-training quantization from one fixed
checkpoint, it reaches \num{0.8304} accuracy against \num{0.8571} in this thesis's
own full-precision baseline. Matching Yuan et al.'s figure is not the target; the
comparison this thesis draws is between its own deployed and full-precision numbers,
on its own device.

The remainder of the thesis is organized as follows. Section~\ref{sec:lit} reviews,
in more detail than this introduction allows, the ferroelectric-memory and
neuromorphic-computing literature that motivates the device and network choices made
here. Section~\ref{sec:method} describes the device structure, the TCAD simulation
setup, the calibration procedure and the extraction rules used to turn simulator
output into the numbers quoted throughout. Section~\ref{sec:device} reports the
device's DC characteristics, its write energy and its analog behaviour, including
the eight-versus-sixteen level result. Section~\ref{sec:network} deploys the
measured device onto the spiking classifier and reports the three results summarized
above, together with a full system energy budget. Section~\ref{sec:discussion} draws
out what a process engineer and a system designer should take from those results and
states plainly the limitations of a two-dimensional, Preisach-model, single-workload
study (Section~\ref{sec:limits}). Section~\ref{sec:conclusion} closes the thesis.
